\documentclass[10pt,letterpaper]{article}
\usepackage[margin=0.67in]{geometry}
\usepackage{booktabs,amsmath,microtype}
\usepackage[hidelinks]{hyperref}
\usepackage{titlesec}
\titleformat{\section}{\normalsize\bfseries}{}{0pt}{}
\titlespacing*{\section}{0pt}{6pt}{3pt}
\setlength{\parindent}{0pt}
\setlength{\parskip}{3pt}
\pagestyle{empty}
\begin{document}
\begin{center}
{\large\bfseries Probing Clinical Meaning in Pretrained Language Representations:\\
Negation, Uncertainty, and Terminology Generalization}\\[3pt]
{\small CSI 436/536 Project P \quad Midterm Report \quad October 20, 2026}\\
{\small Group ID: \textbf{8} \quad Members: \textbf{Yuhan Sun, Zeyu Pan, Jiayu Liao}}
\end{center}
\vspace{-6pt}
\section{Problem Statement}
We investigate whether frozen pretrained representations preserve useful information
about negation and uncertainty, improve learning with limited labeled training data,
and eventually support terminology generalization. The implemented pilot predicts
\emph{sentence-level annotated cue presence} in BioScope abstracts [1]:
\texttt{no\_cue}, \texttt{negation\_only}, or \texttt{speculation\_only}.
These labels do not establish a particular patient's disease status. Target-specific
assertion and unseen-disease evaluation require additional entity-level labels.
\section{Data and Proposed Method}
We downloaded the official BioScope XML and parsed 11,871 sentences from 1,273 abstracts,
including titles. Removing 295 mixed-cue sentences and 17 redundant duplicates leaves
11,559 sentences. Documents connected by duplicate sentences stay in one partition.
Five fixed stratified group folds allocate three folds to training, one to validation,
and one to test. The pilot samples 1,500/400/600 sentences within those partitions,
approximately preserving class proportions. The test contains 439 no-cue, 68 negation,
and 93 speculation sentences.

We compare word unigram/bigram TF-IDF with frozen BERT-base-cased [2], BioBERT [3],
and PubMedBERT [4] features. The encoders use final-layer mean pooling,
excluding padding and special tokens, with a 192-token limit. We fit feature
preprocessing only on each labeled training subset. Class-balanced logistic regression
(LR), plus a TF-IDF linear SVM, select $C\in\{0.1,1,10\}$ on validation macro-F1.
\section{Current Progress and Verified Pilot Results}
The Python pipeline, embedding caches, pinned model revisions, split manifests,
prediction files, and automatic figures are implemented. All five methods completed
10/25/50/100 examples per class and the full pilot training pool with three sampling
seeds: 75 evaluations in total. Integrity checks confirm shared training IDs and
recompute metrics from saved predictions.
\begin{center}
\small
\begin{tabular}{lcc}
\toprule
Method & 100/class: macro-F1 (\%) & Full pool: macro-F1 (\%)\\
\midrule
TF-IDF + LR & $59.1\pm0.8$ & 81.7 \\
TF-IDF + SVM & $59.0\pm0.8$ & 84.3 \\
BERT + LR & $64.4\pm1.3$ & 77.7 \\
BioBERT + LR & $67.8\pm1.1$ & 80.4 \\
PubMedBERT + LR & $65.8\pm1.7$ & 82.6 \\
\bottomrule
\end{tabular}
\end{center}
\vspace{-6pt}
{\small Values at 100/class are mean $\pm$ sample SD over three training subsets on
one fixed test set. Full pool means 1,500 training sentences with natural imbalance;
its repetitions are deterministic. All runs also use 400 labeled validation sentences,
so the budgets count training labels only.}

The pilot suggests that the relative benefit depends on the training setting.
The full-pool comparison also changes class composition, so its difference from
balanced-budget runs is not a pure sample-size effect. These preliminary results
measure cue detection and do not demonstrate clinical understanding.
\section{Challenges and Next Steps}
Explicit cue words may permit lexical shortcuts. Mixed cues, class imbalance,
scope ambiguity, and possible overlap with biomedical pretraining text limit interpretation.
Exact duplicate removal does not guarantee removal of all near-duplicate templates.
Next we will run the uncapped corpus, retain mixed labels in a multilabel task,
and analyze errors by cue family. Ten authored minimal pairs support qualitative
diagnostics only. Disease-held-out tests and target pooling need suitable annotated
targets. We will preserve document and template separation for those extensions.
\vspace{3pt}
\hrule
\vspace{3pt}
{\fontsize{8}{9}\selectfont
[1] Vincze et al., BioScope, \emph{BMC Bioinformatics}, 2008,
\href{https://doi.org/10.1186/1471-2105-9-S11-S9}{doi:10.1186/1471-2105-9-S11-S9}.
[2] Devlin et al., BERT, NAACL 2019.
[3] Lee et al., BioBERT, \emph{Bioinformatics}, 2020.
[4] Gu et al., Domain-Specific Language Model Pretraining, \emph{ACM TCH}, 2021.
Code authored with Codex assistance; group review and course-policy compliance remain
the authors' responsibility.}
\end{document}
